\documentclass[11pt]{article}

\usepackage{amsmath,amssymb}
\usepackage{xurl}
\usepackage[hidelinks]{hyperref}
\setlength{\emergencystretch}{2em}

% Shared commands for notes in docs/paper-gaps/.
% Notation follows the LDT paper (references/ldt-paper/).

\newcommand{\Lean}{\textsc{Lean}}
\newcommand{\leanid}[1]{\textsf{#1}}
\newcommand{\ghissue}[1]{\href{https://github.com/LionSR/MIPStarRE/issues/#1}{GitHub issue \##1}}
\newcommand{\ghpr}[1]{\href{https://github.com/LionSR/MIPStarRE/pull/#1}{GitHub PR \##1}}

% --- Paper-compatible notation ---
\newcommand{\F}{\mathbb{F}}
\newcommand{\N}{\mathbb{N}}
\newcommand{\C}{\mathbb{C}}
\newcommand{\tr}{\operatorname{tr}}
\newcommand{\ot}{\otimes}
\newcommand{\E}{\mathop{\bf E\/}}
\newcommand{\bx}{\mathbf{x}}
\newcommand{\by}{\mathbf{y}}
\newcommand{\bu}{\mathbf{u}}
\newcommand{\bv}{\mathbf{v}}

% Approximation relation (paper uses \approx_\eps and \simeq_\zeta)
\newcommand{\approxeps}{\approx_{\varepsilon}}
\newcommand{\simgeq}{\simeq_{\zeta}}

% Polynomial spaces (paper uses \polyfunc, \polysub, \polymeas)
\newcommand{\polyfunc}[3]{\operatorname{Poly}(#1,#2,#3)}
\newcommand{\polysub}[3]{\operatorname{PolySub}(#1,#2,#3)}
\newcommand{\polymeas}[3]{\operatorname{PolyMeas}(#1,#2,#3)}


\title{The Accumulated Error in the Pasting From-\texorpdfstring{$H$}{H}-to-\texorpdfstring{$G$}{G} Lemma}
\author{}
\date{2026-05-01}

\begin{document}
\maketitle

\section{The assertion}

This note concerns the from-$H$-to-$G$ lemma in the pasting section of
\cite{Ji2020LowIndividualDegree}.%
\footnote{Tracked as part of the discrepancy audit in \ghissue{930}. In the
TeX source consulted here, the statement of the lemma is in
\path{references/ldt-paper/ld-pasting.tex}, lines~1294--1308, and the relevant
error calculation is in lines~1370--1376.}
Here $m$ is the number of variables in each slice, $q$ is the field size,
$d$ is the individual-degree bound, and $k$ is the number of pasting stages.
The parameters $\gamma$ and $\zeta$ are the error parameters carried into the
commutation and self-consistency estimates used by the pasting argument. The operator $G$ is the average slice
submeasurement, and $\widehat H$ denotes the sandwiched completed-slice product
summed over all completed outcomes.

The lemma compares the all-outcomes expansion of the pasted submeasurement with
a Bernoulli polynomial in $G$. It states the error
parameter
\[
  \nu_8
  = 46km(\gamma^{1/32}+\zeta^{1/32}+(d/q)^{1/32}).
  \tag{\path{lem:from-H-to-G}}
\]
The proof first proves an adjacent-stage estimate with error
$2\sqrt{2\zeta}+2\sqrt{\nu_4}$.  Here $\nu_4$ is the commutation error from the
preceding half-sandwich lemma:
\[
  \nu_4
  =426k^2m(\gamma^{1/16}+\zeta^{1/16}+(d/q)^{1/16}).
\]
It then applies the adjacent estimate for $k$ stages.

\section{The arithmetic}

The accumulated adjacent-stage error is
\[
  k(2\sqrt{2\zeta}+2\sqrt{\nu_4}).
\]
Substituting the displayed value of $\nu_4$ gives
\[
\begin{aligned}
  k(2\sqrt{2\zeta}+2\sqrt{\nu_4})
  &= 2k\sqrt{2\zeta}
     +2k\sqrt{426k^2m
       (\gamma^{1/16}+\zeta^{1/16}+(d/q)^{1/16})} \\
  &\le 4k\zeta^{1/32}
     +42k^2m(\gamma^{1/32}+\zeta^{1/32}+(d/q)^{1/32}) \\
  &\le 46k^2m(\gamma^{1/32}+\zeta^{1/32}+(d/q)^{1/32}),
\end{aligned}
\]
where we use the explicit assumptions
\[
  0 \le \gamma,\zeta,d/q \le 1
  \qquad\text{and}\qquad
  k \ge 1,\quad m \ge 1.
\]
The bounds $\gamma,\zeta,d/q \le 1$ allow the square-root powers to be
weakened to $1/32$-powers, and $k,m\ge 1$ allow the first term to be absorbed
into the final $46k^2m(\cdots)$ bound. The second term is quadratic in $k$,
because the factor $k$ from the telescope multiplies the $k$ already present
in $\sqrt{\nu_4}$. The displayed calculation in the paper drops this outer
factor from the commutation term.

Thus the statement proved by the displayed telescope has the corrected error
\[
  \nu_8'
  = 46k^2m(\gamma^{1/32}+\zeta^{1/32}+(d/q)^{1/32}).
\]
This correction does not affect the final pasting theorem: the surrounding
completeness argument absorbs the from-$H$-to-$G$ loss together with the
all-outcomes loss into
\[
  \nu
  =100k^2m\bigl(\eps^{1/32}+\delta^{1/32}+\gamma^{1/32}
      +\zeta^{1/32}+(d/q)^{1/32}\bigr).
\]
The corrected $\nu_8'$ has the same quadratic dependence on $k$ as this
final parameter.

\section{The formal and blueprint statements}

The blueprint records the corrected quadratic error in the from-$H$-to-$G$
lemma and explicitly notes the arithmetic correction.%
\footnote{See \path{blueprint/src/chapter/ch09_pasting.tex}, lines~1037--1100.}
The formal development uses the same corrected quantity.%
\footnote{The corrected error is \leanid{MIPStarRE.LDT.Pasting.fromHToGError}
in \doclink{MIPStarRE/LDT/Pasting/Statements.lean}, lines~113--124 of the
source file. The public lemma \leanid{MIPStarRE.LDT.Pasting.fromHToG} proves
the scalar comparison with this error in
\doclink{MIPStarRE/LDT/Pasting/Bernoulli/FromHToG.lean}. The absorption into
$\nu$ is carried by the final completeness proof in
\doclink{MIPStarRE/LDT/Pasting/Bernoulli/Final.lean}.}
It separates the literal telescoped error from the displayed bound: the
telescope contributes $k$ adjacent-stage losses, and the final comparison
bounds that telescope by the quadratic $\nu_8'$.

\section{Conclusion}

The discrepancy is a harmless intermediate arithmetic correction. The
linear-in-$k$ value of $\nu_8$ stated in the paper does not follow from the
preceding adjacent-stage estimate, because the commutation error $\nu_4$
already contains $k^2$. The corrected bound is quadratic in $k$. Since the
main pasting theorem already uses a quadratic-in-$k$ error parameter with
sufficient slack, the final theorem and the formal statement remain aligned.

\bibliographystyle{alpha}
\bibliography{references}

\end{document}